%% ============================================================
%% Informe Práctica 04 — Dog Messenger
%% CC4P1 Programación Concurrente y Distribuida  2026-I
%% Escuela de Ciencias de la Computación — UNI
%% ============================================================
\documentclass[12pt,a4paper]{article}

\usepackage[spanish,es-noshorthands]{babel}
\usepackage[utf8]{inputenc}
\usepackage[T1]{fontenc}
\usepackage{geometry}
\usepackage{listings}
\usepackage{xcolor}
\usepackage{hyperref}
\usepackage{graphicx}
\usepackage{booktabs}
\usepackage{array}
\usepackage{enumitem}
\usepackage{fancyhdr}
\usepackage{titlesec}
\usepackage{amsmath}
\usepackage{amssymb}
\usepackage{float}
\usepackage{caption}
\usepackage{subcaption}
\usepackage{mdframed}
\usepackage{tikz}
\usetikzlibrary{shapes.geometric, arrows.meta, positioning}

\geometry{
  top=2.5cm, bottom=2.5cm,
  left=2.8cm, right=2.8cm,
  headheight=14pt
}

%% ── Colores ──────────────────────────────────────────────────────────────────
\definecolor{javabg}{RGB}{248,248,255}
\definecolor{kotlinbg}{RGB}{248,255,248}
\definecolor{pybg}{RGB}{255,252,240}
\definecolor{commentcolor}{RGB}{100,120,100}
\definecolor{keywordcolor}{RGB}{0,0,180}
\definecolor{stringcolor}{RGB}{160,32,32}
\definecolor{framecolor}{RGB}{180,180,200}
\definecolor{accentblue}{RGB}{30,90,200}
\definecolor{accentgreen}{RGB}{20,140,60}
\definecolor{accentorange}{RGB}{200,100,0}

%% ── Estilos listings ─────────────────────────────────────────────────────────
\lstdefinestyle{java}{
  language=Java,
  backgroundcolor=\color{javabg},
  basicstyle=\ttfamily\footnotesize,
  keywordstyle=\color{keywordcolor}\bfseries,
  commentstyle=\color{commentcolor}\itshape,
  stringstyle=\color{stringcolor},
  numberstyle=\tiny\color{gray},
  numbers=left, stepnumber=1,
  breaklines=true, breakatwhitespace=true,
  frame=single, rulecolor=\color{framecolor},
  tabsize=4, showstringspaces=false,
  captionpos=b,
  morekeywords={var,val,record,sealed,permits,yield}
}

\lstdefinestyle{kotlin}{
  language=Java,
  backgroundcolor=\color{kotlinbg},
  basicstyle=\ttfamily\footnotesize,
  keywordstyle=\color{accentgreen}\bfseries,
  commentstyle=\color{commentcolor}\itshape,
  stringstyle=\color{stringcolor},
  numberstyle=\tiny\color{gray},
  numbers=left, stepnumber=1,
  breaklines=true, breakatwhitespace=true,
  frame=single, rulecolor=\color{framecolor},
  tabsize=4, showstringspaces=false,
  captionpos=b,
  morekeywords={fun,val,var,data,class,object,companion,override,suspend,launch,coroutineScope}
}

\lstdefinestyle{python}{
  language=Python,
  backgroundcolor=\color{pybg},
  basicstyle=\ttfamily\footnotesize,
  keywordstyle=\color{accentorange}\bfseries,
  commentstyle=\color{commentcolor}\itshape,
  stringstyle=\color{stringcolor},
  numberstyle=\tiny\color{gray},
  numbers=left, stepnumber=1,
  breaklines=true, breakatwhitespace=true,
  frame=single, rulecolor=\color{framecolor},
  tabsize=4, showstringspaces=false,
  captionpos=b
}

%% ── Soporte UTF-8 en listings (acentos españoles) ────────────────────────────
\lstset{
  literate=
    {á}{{\'{a}}}1 {é}{{\'{e}}}1 {í}{{\'{i}}}1 {ó}{{\'{o}}}1 {ú}{{\'{u}}}1
    {Á}{{\'{A}}}1 {É}{{\'{E}}}1 {Í}{{\'{I}}}1 {Ó}{{\'{O}}}1 {Ú}{{\'{U}}}1
    {ñ}{{\~{n}}}1 {Ñ}{{\~{N}}}1 {ü}{{\"u}}1    {Ü}{{\"U}}1
}

%% ── Cabecera / Pie ───────────────────────────────────────────────────────────
\pagestyle{fancy}
\fancyhf{}
\fancyhead[L]{\small CC4P1 --- Dog Messenger}
\fancyhead[R]{\small Práctica 04 · 2026-I}
\fancyfoot[C]{\thepage}
\renewcommand{\headrulewidth}{0.4pt}

%% ── Secciones ────────────────────────────────────────────────────────────────
\titleformat{\section}{\large\bfseries\color{accentblue}}{}{0em}
  {\thesection\enspace}[\titlerule]
\titleformat{\subsection}{\normalsize\bfseries\color{accentblue!80}}{}{0em}
  {\thesubsection\enspace}

%% ── Recuadro de nota ─────────────────────────────────────────────────────────
\newmdenv[
  backgroundcolor=blue!5,
  linecolor=accentblue,
  linewidth=1.2pt,
  innerleftmargin=10pt,
  innerrightmargin=10pt,
  innertopmargin=6pt,
  innerbottommargin=6pt,
  roundcorner=4pt
]{notebox}

%% =============================================================================
\begin{document}

%% ── Portada ──────────────────────────────────────────────────────────────────
\hypersetup{pageanchor=false}
\begin{titlepage}
\centering
\vspace*{1cm}

{\Large \textbf{Universidad Nacional de Ingeniería}}\\[0.3cm]
{\large Escuela de Ciencias de la Computación}\\[0.2cm]
{\large CC4P1 --- Programación Concurrente y Distribuida}\\[0.2cm]
{\large 2026-I}

\vspace{2cm}
\rule{\linewidth}{1.5pt}\\[0.4cm]
{\LARGE \textbf{Dog Messenger}}\\[0.3cm]
{\large Sistema de Mensajería Distribuida Desktop y Móvil}\\[0.2cm]
{\large con Nodo de Ventas Automatizado}
\rule{\linewidth}{1.5pt}

\vspace{2cm}

\begin{tabular}{ll}
\textbf{Práctica:}    & 04 \\[4pt]
\textbf{Integrantes:} & Walter Poma Navarro \\
                      & Jesus Santa Cruz Basilio\\
                      & Max Serrano Arostegui\\[4pt]
\textbf{Lenguajes:}   & LP1 --- Java 17 (servidor + cliente desktop) \\
                      & LP2 --- Kotlin / Android (cliente móvil) \\
                      & LP3 --- Python 3.10+ (nodo de ventas) \\[4pt]
\textbf{Fecha:}       & Junio 2026 \\
\end{tabular}

\vspace{2cm}
{\large \textbf{Resumen}}\\[0.3cm]
\begin{minipage}{0.85\linewidth}
\small
Se implementó \emph{Dog Messenger}, un sistema de mensajería en tiempo real
para escritorio y dispositivos móviles que soporta texto, imágenes y archivos.
El sistema emplea sockets TCP crudos, sin frameworks de comunicación,
haciendo uso extensivo de concurrencia mediante hilos para garantizar
alto rendimiento y ausencia de condiciones de carrera.
Se incorporan grupos de chat, clonación por código QR,
encriptación punto a punto (ECDH + AES-256), métricas de rendimiento
y un nodo de ventas automatizado con procesamiento de lenguaje natural
basado en similitud coseno.
\end{minipage}

\vfill

\end{titlepage}

\hypersetup{pageanchor=true}
\tableofcontents
\newpage

%% =============================================================================
\section{Introducción}

Dog Messenger es un sistema de mensajería distribuida desarrollado como
práctica integradora del curso CC4P1. El proyecto demuestra la aplicación
de técnicas de programación concurrente en un escenario real: múltiples
clientes heterogéneos (Java desktop, Android/Kotlin, Python bot) conectados
a un servidor central a través de sockets TCP binarios.

El principal desafío técnico consiste en garantizar seguridad de hilos
cuando varios clientes leen y escriben simultáneamente estructuras de datos
compartidas, sin usar mecanismos de alto nivel como frameworks de mensajería.
El proyecto cumple la restricción del curso: toda comunicación usa sockets
crudos del sistema operativo, sin WebSocket, Socket.io, RabbitMQ ni similares.

\subsection{Funcionalidades implementadas}

\begin{itemize}[leftmargin=1.5cm]
  \item[\checkmark] Chat en tiempo real entre clientes desktop y móvil.
  \item[\checkmark] Transferencia de archivos e imágenes con verificación SHA-256.
  \item[\checkmark] Grupos de chat con múltiples participantes.
  \item[\checkmark] Clonación de historial por código QR entre dispositivos.
  \item[\checkmark] Encriptación punto a punto (ECDH secp256r1 + AES-256/CBC).
  \item[\checkmark] Nodo de ventas autónomo con NLP (similitud coseno, solo stdlib).
  \item[\checkmark] Métricas de rendimiento en tiempo real.
  \item[\checkmark] Despliegue en red LAN con múltiples nodos.
\end{itemize}

%% =============================================================================
\section{Arquitectura del Sistema}

\subsection{Visión general}

La arquitectura sigue el patrón cliente-servidor centralizado (estrella).
El servidor Java actúa como único punto de enrutamiento: los clientes no se
comunican directamente entre sí, sino que envían mensajes al servidor y este
los reenvía al destinatario. Esta decisión simplifica la gestión de conexiones
y permite que el servidor contabilice métricas globales sin requerir
coordinación entre clientes.

\begin{figure}[H]
\centering
\begin{tikzpicture}[
  node distance=1.8cm and 2.2cm,
  box/.style={rectangle, rounded corners=4pt, draw, minimum width=2.6cm,
              minimum height=1cm, align=center, font=\small},
  server/.style={box, fill=blue!15, draw=accentblue, thick},
  client/.style={box, fill=green!10, draw=accentgreen},
  bot/.style={box, fill=orange!15, draw=accentorange},
  arr/.style={-Latex, thick}
]

\node[server] (srv) {\textbf{Servidor}\\Java 17};

\node[client, above left=of srv]  (dc1) {Desktop\\Java};
\node[client, above right=of srv] (dc2) {Desktop\\Java};
\node[client, below left=of srv]  (and) {Android\\Kotlin};
\node[bot,    below right=of srv] (bot) {Nodo Ventas\\Python};

\draw[arr,<->] (srv) -- (dc1) node[midway,above left,font=\tiny]{TCP};
\draw[arr,<->] (srv) -- (dc2) node[midway,above right,font=\tiny]{TCP};
\draw[arr,<->] (srv) -- (and) node[midway,below left,font=\tiny]{TCP};
\draw[arr,<->] (srv) -- (bot) node[midway,below right,font=\tiny]{TCP};

\end{tikzpicture}
\caption{Arquitectura estrella: todos los nodos se conectan al servidor central.}
\end{figure}

\subsection{Patrones de concurrencia utilizados}

\begin{description}[leftmargin=2.5cm,labelwidth=2.5cm]
  \item[Thread-per-Connection] El servidor crea un \texttt{ClientHandler} que
    extiende \texttt{Thread} por cada cliente aceptado. Permite atender
    múltiples clientes en paralelo sin bloquear el hilo de aceptación.
  \item[Producer-Consumer] El cliente desktop utiliza \texttt{SenderThread}
    con una \texttt{BlockingQueue}. El hilo de la interfaz gráfica produce
    mensajes; el hilo dedicado de envío los consume de forma ordenada, sin
    condiciones de carrera sobre el socket.
  \item[Thread-safe Registry] \texttt{ConcurrentHashMap} y \texttt{AtomicInteger}
    permiten que múltiples \texttt{ClientHandler} registren y consulten usuarios
    simultáneamente sin bloqueos explícitos.
\end{description}

%% =============================================================================
\section{Protocolo de Comunicación Binario}

\subsection{Formato del header}

Todos los mensajes usan un header binario de 9 bytes en orden big-endian.
Se eligió un protocolo binario en lugar de texto (JSON/XML) porque reduce
el overhead de procesamiento, elimina ambigüedades de codificación y es
compatible de forma nativa con \texttt{DataInputStream} / \texttt{DataOutputStream}
tanto en Java como en Kotlin. En Python, el módulo \texttt{struct} permite
leer y escribir el mismo formato con \texttt{struct.pack('>iBHH', ...)}.

\begin{center}
\begin{tabular}{|c|c|c|c|}
\hline
\textbf{payload\_len} & \textbf{type} & \textbf{sender\_id} & \textbf{receiver\_id} \\
4 bytes (int) & 1 byte & 2 bytes (short) & 2 bytes (short) \\
\hline
\end{tabular}
\end{center}

\subsection{Tipos de mensaje}

\begin{center}
\begin{tabular}{cll}
\toprule
\textbf{Código} & \textbf{Nombre} & \textbf{Descripción} \\
\midrule
1  & AUTH         & Autenticación (usuario hacia servidor) \\
2  & TEXT         & Mensaje de texto unicast o broadcast \\
3  & FILE\_START  & Inicio de transferencia (nombre + tamaño) \\
4  & FILE\_CHUNK  & Fragmento de archivo (máx. 4096 bytes) \\
5  & FILE\_END    & Fin de transferencia + hash SHA-256 \\
6  & GROUP        & Operaciones de grupo (sub-comandos 0x01--0x04) \\
7  & QR           & Generación y canje de token QR \\
8  & SALES        & Comandos al nodo de ventas \\
9  & METRICS      & Solicitud de métricas al servidor \\
10 & ERROR        & Mensaje de error del servidor \\
11 & DH\_EXCHANGE & Intercambio de clave pública Diffie-Hellman \\
\bottomrule
\end{tabular}
\end{center}

\subsection{Implementación del parser}

El \texttt{MessageParser} encapsula la lectura y escritura del protocolo.
El método \texttt{readHeader} usa \texttt{DataInputStream.readInt()} en lugar
de un \texttt{read(4)} manual porque TCP es un protocolo de flujo: un solo
\texttt{read()} puede devolver menos bytes de los solicitados si el paquete
llega fragmentado. \texttt{readInt()} llama internamente a \texttt{readFully()},
garantizando que los 4 bytes del campo lleguen completos antes de continuar.

\begin{lstlisting}[style=java, caption={MessageParser.java --- lectura y escritura del protocolo binario}]
public class MessageParser {

    // Lee exactamente el header de 9 bytes
    public static Message readHeader(DataInputStream in) throws IOException {
        int   payloadLen  = in.readInt();          // 4 bytes
        int   typeCode    = in.readUnsignedByte();  // 1 byte
        short senderId    = in.readShort();          // 2 bytes
        short receiverId  = in.readShort();          // 2 bytes
        return new Message(payloadLen,
                           MessageType.fromCode(typeCode),
                           senderId, receiverId);
    }

    // Escribe header + payload de forma atomica
    public static void writeMessage(DataOutputStream out,
                                    Message header,
                                    byte[] payload) throws IOException {
        int len = (payload != null) ? payload.length : 0;
        out.writeInt(len);
        out.writeByte(header.type().getCode());
        out.writeShort(header.senderId());
        out.writeShort(header.receiverId());
        if (payload != null && len > 0) out.write(payload);
        out.flush();
    }
}
\end{lstlisting}

%% =============================================================================
\section{Servidor}

\subsection{Bucle de aceptación}

El servidor utiliza el patrón Thread-per-Connection: por cada socket aceptado
se construye un \texttt{ClientHandler} que extiende \texttt{Thread} y se
inicia de forma independiente. El hilo principal solo acepta conexiones y
delega todo el procesamiento a los hilos hijos, lo que permite atender
decenas de clientes simultáneos sin que ninguno bloquee a otro.

\begin{lstlisting}[style=java, caption={Server.java --- bucle de aceptación (Thread-per-Connection)}]
try (ServerSocket server = new ServerSocket(port)) {
    System.out.println("[Server] Escuchando en puerto " + port);
    while (true) {
        Socket client = server.accept();          // bloquea hasta nueva conexion
        MetricsCollector.getInstance().recordConnection();
        new ClientHandler(client).start();        // nuevo hilo por cliente
    }
}
\end{lstlisting}

\subsection{Registro de usuarios con estructuras concurrentes}

El registro de usuarios utiliza \texttt{ConcurrentHashMap} en lugar de un
\texttt{HashMap} con \texttt{synchronized} global, porque permite operaciones
de lectura y escritura concurrentes sobre distintas entradas sin bloquear toda
la tabla. La generación de IDs usa \texttt{AtomicInteger.getAndIncrement()},
que es atómica a nivel de hardware y evita la posibilidad de que dos clientes
que se conectan al mismo tiempo reciban el mismo ID.

\begin{lstlisting}[style=java, caption={UserRegistry.java --- ConcurrentHashMap + AtomicInteger}]
public class UserRegistry {
    // ConcurrentHashMap: lecturas/escrituras simultaneas sin sincronizacion global
    private final ConcurrentHashMap<Short, ClientHandler> handlers  =
                                            new ConcurrentHashMap<>();
    private final ConcurrentHashMap<String, Short>        usernames =
                                            new ConcurrentHashMap<>();
    // AtomicInteger: IDs unicos aunque dos clientes se conecten a la vez
    private final AtomicInteger idCounter = new AtomicInteger(1);

    public short register(String username, ClientHandler handler) {
        short newId = (short) idCounter.getAndIncrement(); // atomico
        if (usernames.putIfAbsent(username, newId) != null)
            return -1;  // username ya existe
        handlers.put(newId, handler);
        return newId;
    }
}
\end{lstlisting}

El método \texttt{putIfAbsent} es atómico: aunque dos hilos ejecuten
\texttt{register} con el mismo nombre de usuario simultáneamente, solo uno
logrará insertar el valor y el otro recibirá \texttt{-1}, evitando duplicados
sin necesidad de un bloque \texttt{synchronized}.

\subsection{Dispatch de mensajes entrantes}

Cada \texttt{ClientHandler} ejecuta un bucle de lectura en su propio hilo.
Al recibir un mensaje, el método \texttt{dispatch} lo clasifica por tipo y
lo redirige al handler correspondiente. El envío de vuelta al cliente
(\texttt{send}) está declarado \texttt{synchronized} para evitar que dos
hilos distintos entrelacen bytes sobre el mismo socket de salida.

\begin{lstlisting}[style=java, caption={ClientHandler.java --- dispatch según tipo de mensaje}]
private void dispatch(Message header, byte[] payload) throws IOException {
    MetricsCollector.getInstance().recordMessage(header.payloadLen());
    switch (header.type()) {
        case AUTH        -> handleAuth(payload);
        case TEXT        -> handleText(header, payload);
        case FILE_START,
             FILE_CHUNK,
             FILE_END    -> handleFile(header, payload);
        case GROUP       -> handleGroup(header, payload);
        case QR          -> handleQr(header, payload);
        case SALES       -> handleSales(header, payload);
        case METRICS     -> handleMetrics();
        case DH_EXCHANGE -> handleDhExchange(header, payload);
        default          -> sendError(header.senderId(), 0x01, "tipo desconocido");
    }
}

// Envio thread-safe: synchronized evita que dos hilos entreleacen bytes
public synchronized void send(Message header, byte[] payload) throws IOException {
    MessageParser.writeMessage(out, header, payload);
}
\end{lstlisting}

\subsection{Enrutamiento al nodo de ventas}

El nodo de ventas se registra en el servidor como cualquier cliente con el
nombre de usuario \texttt{"ventas"}. Cuando llega un mensaje de tipo SALES,
el servidor busca el handler de ese usuario y reenvía el mensaje de forma
transparente, preservando el \texttt{sender\_id} original para que el bot
sepa a quién responder.

\begin{lstlisting}[style=java, caption={ClientHandler.java --- handleSales()}]
private void handleSales(Message header, byte[] payload) throws IOException {
    Short ventasId = UserRegistry.getInstance().getIdByUsername("ventas");
    if (ventasId == null) {
        sendError(header.senderId(), 0x05, "Nodo de ventas no conectado");
        return;
    }
    ClientHandler ventasHandler = UserRegistry.getInstance()
                                              .getHandler(ventasId);
    // Reenvio transparente: el nodo ventas ve el sender_id original
    ventasHandler.send(header, payload);
}
\end{lstlisting}

%% =============================================================================
\section{Cliente Desktop}

\subsection{Patrón Producer-Consumer para envío}

El cliente desktop separa la interfaz gráfica del envío al socket mediante
el patrón Producer-Consumer. El hilo de la GUI encola mensajes en una
\texttt{LinkedBlockingQueue} sin bloquearse nunca; el \texttt{SenderThread}
los toma de la cola y los envía en orden. Esto evita que una operación de red
lenta congele la interfaz gráfica y garantiza que los mensajes lleguen en el
orden en que fueron generados.

\begin{lstlisting}[style=java, caption={SenderThread.java --- cola de mensajes thread-safe}]
public class SenderThread extends Thread {
    // BlockingQueue: el hilo de UI nunca bloquea al encolar
    private final BlockingQueue<QueuedMessage> queue =
                                    new LinkedBlockingQueue<>();

    public void queueMessage(Message header, byte[] payload) {
        queue.offer(new QueuedMessage(header, payload)); // no bloquea
    }

    @Override
    public void run() {
        while (!shutdown) {
            QueuedMessage msg = queue.take(); // bloquea hasta haber trabajo
            byte[] data = msg.payload();
            // Encripta si hay clave E2E establecida para ese destinatario
            if (crypto != null && crypto.hasKeyFor(msg.header().receiverId())) {
                data = crypto.encrypt(msg.header().receiverId(), data);
            }
            MessageParser.writeMessage(out, msg.header(), data);
        }
    }
}
\end{lstlisting}

\subsection{Encriptación punto a punto}

La encriptación extremo a extremo funciona en dos fases. Primero, los dos
clientes intercambian sus claves públicas ECDH a través del servidor usando
el tipo \texttt{DH\_EXCHANGE}; cada uno deriva de forma independiente el mismo
secreto compartido (propiedad de Diffie-Hellman). Segundo, ese secreto se
procesa con SHA-256 para obtener una clave AES-256 de 32 bytes, con la que se
cifran los mensajes siguientes en modo CBC con IV aleatorio. El servidor solo
ve bytes cifrados y no puede leer el contenido.

\begin{lstlisting}[style=java, caption={CryptoManager.java --- intercambio de claves y cifrado}]
public class CryptoManager {
    private final KeyPair myKeyPair;
    private final Map<Short, SecretKey> sharedKeys = new ConcurrentHashMap<>();

    public CryptoManager() throws GeneralSecurityException {
        KeyPairGenerator kpg = KeyPairGenerator.getInstance("EC");
        kpg.initialize(new ECGenParameterSpec("secp256r1")); // curva NIST P-256
        myKeyPair = kpg.generateKeyPair();
    }

    // Calcula secreto compartido con Diffie-Hellman sobre curva eliptica
    public void computeSharedKey(short peerId, byte[] peerPublicKeyBytes)
            throws GeneralSecurityException {
        KeyFactory kf = KeyFactory.getInstance("EC");
        PublicKey peerPub = kf.generatePublic(
                new X509EncodedKeySpec(peerPublicKeyBytes));
        KeyAgreement ka = KeyAgreement.getInstance("ECDH");
        ka.init(myKeyPair.getPrivate());
        ka.doPhase(peerPub, true);
        byte[] secret = ka.generateSecret();
        // Deriva clave AES-256 del secreto compartido
        byte[] keyBytes = MessageDigest.getInstance("SHA-256").digest(secret);
        sharedKeys.put(peerId, new SecretKeySpec(keyBytes, "AES"));
    }

    // Cifra con AES-256/CBC; payload = IV(16B) || ciphertext
    public byte[] encrypt(short peerId, byte[] plaintext) throws Exception {
        SecretKey key = sharedKeys.get(peerId);
        Cipher cipher = Cipher.getInstance("AES/CBC/PKCS5Padding");
        cipher.init(Cipher.ENCRYPT_MODE, key);  // genera IV aleatorio
        byte[] iv  = cipher.getIV();
        byte[] ct  = cipher.doFinal(plaintext);
        byte[] out = new byte[16 + ct.length];
        System.arraycopy(iv, 0, out, 0, 16);
        System.arraycopy(ct, 0, out, 16, ct.length);
        return out;
    }
}
\end{lstlisting}

\subsection{Visualización de imágenes inline}

El log de chat usa \texttt{JTextPane} en lugar de \texttt{JTextArea} porque
\texttt{JTextPane} permite insertar componentes Swing arbitrarios mediante
\texttt{insertComponent()}. Cuando llega la transferencia completa de un
archivo de imagen, el cliente lo detecta por extensión y construye un
\texttt{JLabel} con el \texttt{ImageIcon} escalado a un máximo de 300 px de
ancho, que inserta directamente en el panel de chat.

\begin{lstlisting}[style=java, caption={ChatFrame.java --- detección y renderizado de imágenes recibidas}]
// Al completar FILE_END, comprueba si es imagen por extension
case FILE_END -> {
    byte[] fileBytes = activeDownloadStream.toByteArray();
    Files.write(dest.toPath(), fileBytes);
    if (isImage(activeDownloadFilename)) {
        appendLog("[imagen de #" + h.senderId() + "]");
        appendImage(fileBytes);
    } else {
        appendLog("[archivo] " + activeDownloadFilename + " guardado.");
    }
}

private void appendImage(byte[] imageBytes) {
    ImageIcon raw = new ImageIcon(imageBytes);
    int w = Math.min(raw.getIconWidth(), 300); // maximo 300px de ancho
    int h = (int)(raw.getIconHeight() * ((double) w / raw.getIconWidth()));
    Image scaled = raw.getImage().getScaledInstance(w, h, Image.SCALE_SMOOTH);
    JLabel imgLabel = new JLabel(new ImageIcon(scaled));
    txtChatLog.insertComponent(imgLabel); // JTextPane permite componentes inline
    appendLog("");
}

private static boolean isImage(String filename) {
    String lo = filename.toLowerCase();
    return lo.endsWith(".png") || lo.endsWith(".jpg") || lo.endsWith(".jpeg")
        || lo.endsWith(".gif") || lo.endsWith(".bmp");
}
\end{lstlisting}

\subsection{Capturas del software}

A continuación se muestran capturas de pantalla del cliente desktop en
ejecución.

\begin{figure}[H]
\centering
\includegraphics[width=\textwidth]{captura-chat.png}
\caption{Ventana principal de chat con mensajes de texto e imagen inline.}
\end{figure}

\begin{figure}[H]
\centering
\includegraphics[width=\textwidth]{captura-ventas.png}
\caption{Conversación con el nodo de ventas mostrando flujo de pedido completado.}
\end{figure}

%% =============================================================================
\section{Cliente Móvil Android}

\subsection{Compatibilidad de protocolo}

El cliente Android implementa exactamente el mismo protocolo de 9 bytes que
el servidor Java. No se reescribió el protocolo: se usó \texttt{DataInputStream}
y \texttt{DataOutputStream} de la JVM con las mismas llamadas que en Java, lo
que garantiza compatibilidad byte a byte sin conversiones adicionales.

\begin{lstlisting}[style=kotlin, caption={MessageParser.kt --- mismo protocolo en Kotlin}]
object MessageParser {
    fun readHeader(input: DataInputStream): Message {
        val payloadLen  = input.readInt()              // 4 bytes
        val typeCode    = input.readByte().toInt() and 0xFF
        val senderId    = input.readShort()
        val receiverId  = input.readShort()
        return Message(payloadLen,
                       MessageType.fromCode(typeCode),
                       senderId, receiverId)
    }

    fun writeMessage(output: DataOutputStream,
                     header: Message,
                     payload: ByteArray?) {
        output.writeInt(payload?.size ?: 0)
        output.writeByte(header.type.code)
        output.writeShort(header.senderId.toInt())
        output.writeShort(header.receiverId.toInt())
        payload?.let { output.write(it) }
        output.flush()
    }
}
\end{lstlisting}

\subsection{Envío thread-safe en Android}

Android no tiene \texttt{BlockingQueue} en su API de UI, pero sí permite
usar \texttt{Thread} de Java. El envío de mensajes se protege con
\texttt{synchronized(this)} sobre la conexión, que es suficiente para
garantizar que dos hilos no entrelacen bytes sobre el mismo
\texttt{DataOutputStream}.

\begin{lstlisting}[style=kotlin, caption={Connection.kt --- synchronized para escritura al socket}]
private fun sendMessage(header: Message, payload: ByteArray) {
    Thread {
        synchronized(this) {
            // synchronized garantiza que solo un Thread escribe a la vez
            output?.let { MessageParser.writeMessage(it, header, payload) }
        }
    }.start()
}

// Ejemplo: enviar mensaje de grupo
fun sendGroupMessage(groupId: Short, message: String) {
    val gidBytes = ByteBuffer.allocate(2)
                             .putShort(groupId).array()
    val msgBytes = message.toByteArray(Charsets.UTF_8)
    val payload  = byteArrayOf(0x03) + gidBytes + msgBytes
    sendMessage(Message(payload.size, MessageType.GROUP,
                        userId, 0xFFFF.toShort()), payload)
}
\end{lstlisting}

\subsection{Funcionalidades implementadas en Android}

\begin{itemize}
  \item Autenticación con nombre de usuario y asignación de ID por el servidor.
  \item Chat de texto unicast y broadcast.
  \item Transferencia de archivos con verificación SHA-256.
  \item Grupos: crear, unirse, enviar y listar.
  \item QR: generar token, canjear token de otro dispositivo.
  \item Métricas del servidor.
  \item Guardado de archivos en \texttt{MediaStore} (Android API 30+).
\end{itemize}

%% =============================================================================
\section{Nodo de Ventas}

\subsection{Integración como cliente del servidor}

El nodo de ventas no requiere un protocolo separado ni un puerto distinto.
Se conecta al servidor exactamente igual que cualquier cliente de chat,
se autentica con el nombre de usuario \texttt{"ventas"} y queda registrado
en el \texttt{UserRegistry}. Cualquier usuario del chat puede enviarle
mensajes de texto escribiendo su ID en el campo \emph{Para}. Esta decisión
evita añadir lógica especial al servidor y reutiliza toda la infraestructura
de enrutamiento existente.

\begin{lstlisting}[style=python, caption={main.py --- autenticación y dispatch del nodo}]
class SalesNode:
    def _authenticate(self):
        # Se autentica como cualquier usuario normal
        send_message(self.sock, AUTH, 0, BROADCAST, b'ventas')

        type_code, _, _, resp = read_message(self.sock)
        self.user_id = (resp[1] << 8) | resp[2]
        print(f"[Sales] Autenticado. ID: #{self.user_id}")
        print(f"[Sales] Escribe el ID #{self.user_id} en 'Para' para chatear")

    def _dispatch(self, type_code, sender_id, payload):
        if type_code == TEXT:
            # Mensajes de texto: responde el chatbot NLP
            text     = payload.decode('utf-8', errors='replace')
            response = self.chatbot.handle(sender_id, text)
            self._send_text(sender_id, response)
        elif type_code == SALES:
            # Sub-comandos estructurados del SalesDialog
            self._handle_sales(sender_id, payload)
\end{lstlisting}

\subsection{Motor NLP con similitud coseno}

Para detectar la intención del usuario sin librerías externas se implementó
similitud coseno sobre vectores bag-of-words. Cada mensaje se tokeniza
(eliminación de stopwords, normalización de acentos, conversión de dígitos a
un token \texttt{NUM}), y el vector resultante se compara con un índice
precalculado al arrancar: un \texttt{Counter} agregado de todas las frases
de entrenamiento por intención. La similitud coseno entre dos vectores es
el coseno del ángulo que forman, lo que mide cuánto vocabulario comparten
proporcionalmente. Se eligió este método porque no requiere entrenamiento
iterativo, es determinista y funciona con la stdlib de Python (\texttt{math},
\texttt{re}, \texttt{collections}).

\begin{lstlisting}[style=python, caption={nlp.py --- similitud coseno con bag-of-words}]
import math, re
from collections import Counter
from corpus import STOPWORDS, INTENTS

# Normaliza acentos: "pedicion" -> sin tildes
_ACCENT_MAP = str.maketrans('áéíóúüñÁÉÍÓÚÜÑ', 'aeiouunAEIOUUN')

def tokenize(text: str) -> list:
    text = text.translate(_ACCENT_MAP).lower()
    text = re.sub(r'[^\w\s]', ' ', text)
    text = re.sub(r'\d+', ' NUM ', text)
    return [t for t in text.split()
            if len(t) > 1 and t not in STOPWORDS]

def cosine_sim(v1: Counter, v2: Counter) -> float:
    common = set(v1) & set(v2)
    dot    = sum(v1[k] * v2[k] for k in common)
    mag1   = math.sqrt(sum(v*v for v in v1.values()))
    mag2   = math.sqrt(sum(v*v for v in v2.values()))
    return dot / (mag1 * mag2) if mag1 and mag2 else 0.0

# Indice precalculado al arrancar: {intent: Counter_agregado}
_INTENT_INDEX = {
    intent: sum((Counter(tokenize(p)) for p in phrases), Counter())
    for intent, phrases in INTENTS.items()
}

def detect_intent(text: str, threshold=0.10):
    user_vec = Counter(tokenize(text))
    scores   = [(i, cosine_sim(user_vec, v))
                for i, v in _INTENT_INDEX.items()]
    best     = max(scores, key=lambda x: x[1])
    return best if best[1] >= threshold else ('DESCONOCIDO', 0.0)
\end{lstlisting}

\subsection{Corpus de entrenamiento}

El corpus se orienta a una tienda de mascotas (\emph{Dog Messenger Store})
para coherencia con el nombre del sistema. Incluye 802 stopwords (288
base + 514 del proyecto \texttt{stopwords-iso}) cargadas desde un archivo
de texto plano, sin librerías adicionales. Las palabras clave operativas
como \texttt{pedido} o \texttt{registrar} fueron excluidas explícitamente
de las stopwords para que lleguen al motor NLP sin filtrarse.

\begin{itemize}
  \item 802 stopwords en total.
  \item 324 términos de productos (alimentos, accesorios, higiene, salud).
  \item 172 frases de entrenamiento en 5 intenciones: REGISTRAR, PEDIDO,
    CONSULTAR, REPORTE, AYUDA.
  \item 667 tokens significativos en el índice.
\end{itemize}

\subsection{Máquina de estados conversacional}

El chatbot mantiene el estado de conversación por usuario en un diccionario
protegido por \texttt{threading.Lock}, ya que el nodo puede atender mensajes
de múltiples usuarios concurrentes. Cada usuario tiene su propio objeto
\texttt{\_State} con el paso actual del flujo (registro, pedido, consulta)
y los datos temporales acumulados entre turnos.

El método \texttt{\_idle} evalúa el mensaje en capas: primero palabras clave
exactas que el NLP confundiría (dígitos del menú, términos como
\texttt{"pedido"} que aparecen en varios corpora); luego el motor NLP para
frases largas de lenguaje libre.

\begin{lstlisting}[style=python, caption={chatbot.py --- maquina de estados + NLP}]
class Chatbot:
    def _idle(self, st, text, lo):
        # 1) Greetings directos (vectores de 1 token son muy debiles)
        if set(lo.split()) <= _GREET_TOKENS:
            return MENU_TEXT

        # 2) Opciones numericas del menu
        if lo == '1': st.step = 'REG_NAME';  return REGISTRO_MSG
        if lo in ('2', 'pedido', 'comprar'): st.step = 'ORD_CID'; return PEDIDO_MSG
        if lo in ('3', 'consultar'):         st.step = 'QRY_CID'; return CONSULTA_MSG
        if lo in ('4', 'reporte'):           return self._report()

        # 3) NLP coseno para lenguaje libre (frases largas)
        intent, score = detect_intent(text)
        if   intent == 'REGISTRAR': st.step = 'REG_NAME';  return REGISTRO_MSG
        elif intent == 'PEDIDO':    st.step = 'ORD_CID';   return PEDIDO_MSG
        elif intent == 'CONSULTAR': st.step = 'QRY_CID';   return CONSULTA_MSG
        elif intent == 'REPORTE':   return self._report()
        elif intent == 'AYUDA':     return MENU_TEXT

        return random.choice(RESPONSES_NO_ENTENDIDO)
\end{lstlisting}

\subsection{Ejemplo de conversación}

\begin{center}
\small
\begin{tabular}{p{5.2cm} p{7.4cm}}
\toprule
\textbf{Usuario} & \textbf{Bot (intención detectada)} \\
\midrule
cuales son mis pedidos de musica &
  CONSULTAR (score=0.35) --- \textit{Cual es tu ID de cliente?} \\[4pt]
quiero registrarme &
  REGISTRAR (score=0.62) --- \textit{Cual es tu nombre completo?} \\[4pt]
quisiera comprar croquetas &
  PEDIDO (score=0.50) --- \textit{Tu numero de cliente?} \\[4pt]
dame el reporte de ventas &
  REPORTE (score=0.53) --- \textit{Muestra estadísticas} \\
\bottomrule
\end{tabular}
\end{center}

%% =============================================================================
\section{Métricas de Rendimiento}

El \texttt{MetricsCollector} registra uptime, conexiones totales, mensajes
procesados y bytes transferidos. Usa \texttt{AtomicInteger} y \texttt{AtomicLong}
en lugar de variables protegidas con \texttt{synchronized} porque las
instrucciones CAS (Compare-And-Swap) que utilizan internamente son operaciones
atómicas del procesador: no necesitan adquirir un monitor y por lo tanto no
generan contención cuando muchos hilos actualizan el contador al mismo tiempo.
El cliente obtiene las métricas enviando un mensaje tipo METRICS al servidor.

\begin{lstlisting}[style=java, caption={MetricsCollector.java --- contadores lock-free}]
public class MetricsCollector {
    private final long          startTime        = System.currentTimeMillis();
    private final AtomicInteger totalConnections = new AtomicInteger(0);
    private final AtomicInteger totalMessages    = new AtomicInteger(0);
    private final AtomicLong    totalBytes       = new AtomicLong(0);

    public void recordMessage(int payloadLen) {
        totalMessages.incrementAndGet();
        totalBytes.addAndGet(payloadLen);
    }

    public String toJson() {
        long uptime = (System.currentTimeMillis() - startTime) / 1000;
        return String.format(
            "{\"uptime_s\":%d,\"online\":%d,\"mensajes\":%d,\"bytes\":%d}",
            uptime,
            UserRegistry.getInstance().getOnlineCount(),
            totalMessages.get(),
            totalBytes.get()
        );
    }
}
\end{lstlisting}

%% =============================================================================
\section{Clonación de Chat por QR}

La clonación por QR permite que un usuario transfiera su historial de chat
a otro dispositivo sin necesitar una cuenta externa. El flujo es el siguiente:

\begin{enumerate}
  \item El usuario A solicita un token QR desde el cliente desktop.
  \item El servidor genera un UUID con TTL de 120 segundos y lo almacena
    en \texttt{QrTokenStore}.
  \item El token se envía al cliente A y se renderiza como imagen QR usando
    la librería ZXing.
  \item El usuario B escanea o introduce el token en su cliente.
  \item El servidor valida el token, recupera el historial de A y lo envía a B.
\end{enumerate}

El TTL de 120 segundos se eligió para que el código QR sea válido durante
el tiempo razonable de un escaneo sin que quede abierto indefinidamente.
El token se elimina en cuanto se canjea (\texttt{tokens.remove}) para evitar
reutilización.

\begin{lstlisting}[style=java, caption={QrTokenStore.java --- tokens con TTL de 120s}]
public String generate(short userId) {
    String token = UUID.randomUUID().toString();
    tokens.put(token, new Entry(userId,
                                System.currentTimeMillis() + 120_000));
    return token;
}

public short redeem(String token) {
    Entry entry = tokens.remove(token);
    if (entry == null || System.currentTimeMillis() > entry.expiresAt())
        return -1;   // expirado o no existe
    return entry.userId();
}
\end{lstlisting}

%% =============================================================================
\section{Conclusiones}

Dog Messenger demuestra que es posible construir un sistema de mensajería
distribuido completo usando únicamente los primitivos básicos de cada
lenguaje: sockets crudos e hilos, sin frameworks de comunicación.

La separación entre el protocolo binario y los lenguajes clientes resultó
clave para la interoperabilidad: Java, Kotlin y Python interactúan sin
adaptadores porque todos leen y escriben el mismo header de 9 bytes.
Las estructuras concurrentes de Java (\texttt{ConcurrentHashMap},
\texttt{AtomicInteger}, \texttt{BlockingQueue}) permitieron eliminar
bloqueos globales y garantizar consistencia bajo carga sin sacrificar
rendimiento.

El nodo de ventas en Python ilustra cómo un tercer lenguaje puede integrarse
al sistema sin modificar el servidor, simplemente respetando el contrato
del protocolo.

\end{document}
